\sectionTitle{Work Experience}{\faSuitcase}

\begin{longtable}{R{2.5cm}|p{14.8cm}}

	\emph{Present} & Software Lead --- Diagnostic Systems, \textsc{Hitachi Rail STS S.p.A.}, Naples, Italy\\
	\textsc{2024}  & \emph{Diagnostic systems for metro, tram, and high-speed platforms}\\&
	\footnotesize{
		\begin{itemize}
			\item \textbf{Software lead} for the diagnostic systems of \textbf{Rome Metro (lines A/B)} and \textbf{Florence Tram}, owning the software from development through integration and on-site commissioning
			\item \textbf{Software development support} on the \textbf{Milan Metro (lines 1, 2, 3)} diagnostic system, contributing to development, integration, and issue resolution
			\item Led a focused stabilization phase on the Rome Metro project ahead of official customer testing, \textbf{resolving 230 open issues within one month} and securing the validation milestones on schedule
			\item Led the full software lifecycle of onboard diagnostic systems, developing the HMI and optimizing the embedded backend framework in Embedded C
			\item Integrated the HMI and diagnostic logic for the \textbf{HABD} (Hot Axle Box Detector) and \textbf{DNRA} (Locked Axle Detection) safety systems on the \textbf{ETR1000 Frecciarossa}, covering sensor data handling, alarm management, and driver notification
		\end{itemize}
		\vspace{-1em}
	}\\&
	\footnotesize{\emph{Technologies:} Embedded C, embedded Linux, Ethernet onboard networks, Jira}\\
\multicolumn{2}{c}{}\\

	\textsc{2024} & System Engineer --- Event Recorder \& Safety-Critical Subsystems, \textsc{Hitachi Rail STS S.p.A.}, Naples, Italy\\
	\textsc{2022} & \emph{System ownership of onboard safety-critical products}\\&
	\footnotesize{
		\begin{itemize}
			\item \textbf{System owner} of the \textbf{Event Recorder} product line: defined the complete system architecture and requirements, driving the solution from specification through integration, validation, and customer acceptance for the metro networks of \textbf{Milan, Rome, and Washington (WMATA)}
			\item System Engineer referent for Bogie Control System communication (\textbf{TCMS} integration) on the \textbf{ETR1000 Frecciarossa} high-speed platform, using message data protocol to define the communication architecture and coordinating cross-functional validation across software, hardware, and RAMS teams
			\item Technical referent for onboard safety-critical subsystems --- Event Recorder, Dead Man, and Zero Speed detection --- covering requirements analysis, traceability (DOORS), interface definition, and verification strategy
			\item Coordinated integration, troubleshooting, and commissioning of diagnostic subsystems in operational environments
		\end{itemize}
		\vspace{-1em}
	}\\&
	\footnotesize{\emph{Technologies:} TRDP, IPTCom, TCMS, DOORS, system architecture, IV\&V}\\
\multicolumn{2}{c}{}\\

	\textsc{2022}      & Embedded Software Developer --- Regional \& High-Speed Platforms, \textsc{Hitachi Rail STS S.p.A.}, Naples, Italy\\
	\textsc{Apr. 2018} & \emph{Onboard diagnostic software development}\\&
	\footnotesize{
		\begin{itemize}
			\item Developed and commissioned the diagnostic systems for the regional fleets \textbf{Trenitalia Rock}, \textbf{Trenitalia Blues}, and \textbf{Ferrovie Nord Milano}, and contributed to the high-speed platforms (\emph{ETR1000 Frecciarossa, ETR700})
			\item Developed and extended Condition-Based Maintenance (\textbf{CBM}) software in Embedded C, and designed diagnostic algorithms in C and MATLAB/Simulink
			\item Enhanced and refactored the embedded diagnostic software framework, improving maintainability and reuse across train platforms
			\item Implemented multi-vehicle composition management logic (backend and frontend) and designed diagnostic HMI components in HTML/JavaScript
			\item Improved the reliability and stability of distributed onboard communication networks, managing issue tracking and defect resolution through Jira
		\end{itemize}
		\vspace{-1em}
	}\\&
	\footnotesize{\emph{Technologies:} Embedded C, MATLAB/Simulink, HTML/JavaScript, Jira}\\

\end{longtable}
